\documentclass[11pt,a4paper]{article}
\usepackage[T1]{fontenc}
\usepackage{amsmath,amssymb,geometry}
\geometry{margin=26mm}
\title{EPPA--III, Claim 3.19: the empty/full complementary profiles}
\author{Validation supplement; author manuscript frozen}
\date{16 October 2026}
\begin{document}
\maketitle

\paragraph{Context.}
Let \(H\) be a vertex-transitive graph on ten vertices,
and let the designated induced subgraph \(G\) be \(C_5\).
Write \(B=V(H)\setminus V(G)\) and
\(P_b=N_H(b)\cap V(G)\) for \(b\in B\).
The already verified proper-profile case (PR~\#265)
proves \(H[B]\cong C_5\) if some
\(\varnothing\subsetneq P_b\subsetneq V(G)\).
It remains to treat
\[
       P_b\in\{\varnothing,V(G)\}\qquad(b\in B).
\]
This case needs no EPPA hypothesis: graph regularity is enough.

\section*{Lemma: empty and full profiles cannot coexist}
If \(P_u=V(G)\), then \(u\) has at least five neighbours
in \(G\), so \(\deg_H(u)\ge5\).
If \(P_v=\varnothing\), then all neighbours of \(v\)
lie among the other four vertices of \(B\); hence
\(\deg_H(v)\le4\). This contradicts vertex-transitivity.
Therefore either \emph{every} outside profile is empty
or \emph{every} outside profile is full.

The two degree estimates and the no-mixed-profiles
conclusion are the main targets of
\texttt{EPPAIII/BlockSystems/CycleFiveEmptyFullRigidity.lean}.
They are stated for an arbitrary vertex type \(W\)
with an actual induced copy \(e:\mathrm{Fin}\,5\hookrightarrow W\).
There is no assumed imprimitivity system or
equivariant labelling of \(B\).

\section*{Lemma: the induced outside graph is a five-cycle}
First suppose every \(P_b\) is empty.
Then there are no edges between \(G\) and \(B\).
Since \(G=C_5\) is connected, \(V(G)\) is an entire
connected component of \(H\).
Choose \(x\in V(G)\) and \(b\in B\).
Vertex-transitivity gives \(q\in\operatorname{Aut}(H)\)
with \(q(x)=b\); the image \(q(V(G))\) is the connected
component of \(b\). This component is disjoint from
\(V(G)\), has five vertices, and thus equals \(B\).
Consequently \(H[B]\cong H[G]\cong C_5\).

If every \(P_b\) is full, apply the preceding
argument to the \emph{complement} \(\overline H\),
which is again vertex-transitive and has no cross
edges. The selected induced graph is
\(\overline{C_5}\), which is connected and isomorphic
to \(C_5\). It follows that
\(\overline{H[B]}\cong\overline{C_5}\),
so again \(H[B]\cong C_5\).

\paragraph{Alternative numeric proof.}
In either uniform case, every vertex of \(G\)
has total degree \(2+\varepsilon 5\),
where \(\varepsilon=0\) for empty and \(1\) for full
profiles. Every outside vertex has the same
\(\varepsilon 5\) cross-neighbours.
By regularity, its degree within \(B\) is \(2\).
A simple \(2\)-regular graph on five vertices
is necessarily a single five-cycle
(a disconnected \(2\)-regular graph would have
at least three vertices in each component).
This is an alternative to formalizing component
transport and may yield a shorter Lean proof.

\section*{Consequence for Claim 3.19}
Combining this argument with the already verified
nonempty-proper-profile proof establishes
\[
\boxed{\begin{gathered}
 |V(H)|=10,\quad H\text{ vertex-transitive},\\
 e:C_5\hookrightarrow H\text{ an induced EPPA embedding}
 \quad\Longrightarrow\quad H[B]\cong C_5 .
\end{gathered}}
\]
No imprimitivity hypothesis is needed for
\emph{this complementary-half conclusion}.
It remains essential in the broader host
classification of Lemma~3.18: for instance
the Petersen graph is a primitive possibility.
The argument above is a \emph{paper-level proof};
the full Lean graph-isomorphism conclusion must
not be marked certified until the connected-component
or \(2\)-regular-five-vertex step is formally checked.

\section*{Minimal suggested insertion}
After splitting the outside profiles into proper and
empty/full cases, insert just:
\begin{quote}
If all outside neighbourhood profiles are empty or full,
they must be uniform: a vertex with a full profile has
degree at least five, whereas one with an empty profile
has degree at most four. Thus either \(H\) or
\(\overline H\) has no edges across the two halves.
Transitivity then identifies the two connected
components, so \(H[B]\cong C_5\).
\end{quote}
The established proper-profile proof handles the other
case independently. Leave the rest of the
author's manuscript unchanged.
\section*{Exact Lean assembly target}
The proposed \texttt{CycleFiveEmptyFullConclusion.lean}
module combines
\texttt{CycleFiveEmptyFullRigidity.lean} and
\texttt{TransitiveSeparatedHalf.lean}.
Its final theorem constructs an actual equivalence
\(\mathrm{Fin}\,5\simeq V(H)\setminus V(C_5)\) and proves
graph equality after pullback, rather than asserting a
mere edge-count identity. The complement \(C_5\) is
connected and isomorphic to \(C_5\) by the explicit
multiplication-by-two permutation.
The source-anchored register entry is C23.
No full Lean certification may be recorded until
the combined workflow passes.

\end{document}
